\paragrafos{Άθροισμα και διαφορά τμημάτων: από την ευθεία στην τεθλασμένη}{Β.1.4, «Πρόσθεση και αφαίρεση ευθυγράμμων τμημάτων»}

\begin{anagke}
Ο παππούς Γιάννης φέρνει τα δέντρα που θα φυτευτούν σε μία σειρά, και ο Ηλίας μετράει τις αποστάσεις: το δεύτερο δέντρο απέχει από το πρώτο $3$ μέτρα, το τρίτο $4$ μέτρα από το δεύτερο και το τέταρτο $2{,}5$ από το τρίτο. «Πόσο απέχει το τέταρτο από το πρώτο;» ρωτάει η Πολυξένη. Ο Ηλίας το βρίσκει εύκολα: $3+4+2{,}5=9{,}5$ μέτρα. «Και το τρίτο από το πρώτο; $3+4=7$. Και το τέταρτο από το δεύτερο;». Εδώ ο Στέλιος προσπαθεί να προσθέσει, αλλά η Αγγελική λέει: «Αφαίρεσε: από την απόσταση του τέταρτου από το πρώτο, αφαιρώ την απόσταση του δεύτερου από το πρώτο: $9{,}5-3=6{,}5$».

Η παρέα ανακαλύπτει ότι οι αποστάσεις προσθέτονται και αφαιρούνται όταν τα σημεία βρίσκονται στην ίδια ευθεία με τη σειρά. Μετά ο Νίκος γυρίζει στο μονοπάτι που στρίβει: «Αν πάω από το πρώτο στο τρίτο δέντρο περνώντας από ένα άλλο σημείο που δεν είναι στην ευθεία, κάνω περισσότερο δρόμο;». Η Ρήνη το δοκιμάζει με σχοινί και δείχνει ότι το ευθύγραμμο τμήμα είναι πάντα το συντομότερο μονοπάτι ανάμεσα σε δύο σημεία. Και η κυρία Αλεξάνδρα προσθέτει: «Το μήκος του περιγράμματος ενός παρτεριού, για το φράχτη του, είναι η περίμετρός του».
\end{anagke}

\begin{lexi}{άθροισμα τμημάτων, περίμετρος}{άθροισμα τμημάτων, περίμετρος}
Η \lexiword{περίμετρος} είναι το «γύρω μέτρημα» (περί + μέτρο): το συνολικό μήκος του περιγράμματος ενός σχήματος. Η \lexiword{τεθλασμένη} γραμμή έχει μήκος το άθροισμα των μηκών των ευθυγράμμων τμημάτων της.
\end{lexi}

\begin{metafrasi}
Τα $\mathrm{Α},\,\mathrm{Β},\,\Gamma$ στη σειρά σε ευθεία. & $\mathrm{ΑΓ}=\mathrm{ΑΒ}+\mathrm{ΒΓ}$\\
Ό,τι περισσεύει μετά το $\mathrm{ΑΒ}$. & $\mathrm{ΒΓ}=\mathrm{ΑΓ}-\mathrm{ΑΒ}$\\
Τεθλασμένη $\mathrm{ΑΒΓΔ}$. & μήκος $=\mathrm{ΑΒ}+\mathrm{ΒΓ}+\Gamma\Delta$\\
Ο συντομότερος δρόμος. & $\mathrm{ΑΒ}<$ μήκος κάθε τεθλασμένης με άκρα $\mathrm{Α},\,\mathrm{Β}$\\
Το άθροισμα των πλευρών. & περίμετρος\\
\end{metafrasi}

\tmima{Η θεωρία}

Αν τα σημεία $\mathrm{Α},\,\mathrm{Β},\,\Gamma$ βρίσκονται με αυτή τη σειρά σε μια ευθεία, το τμήμα $\mathrm{ΑΓ}$ λέγεται \textbf{άθροισμα} των τμημάτων $\mathrm{ΑΒ}$ και $\mathrm{ΒΓ}$:
$$\mathrm{ΑΓ}=\mathrm{ΑΒ}+\mathrm{ΒΓ}.$$
Και το $\mathrm{ΒΓ}$ είναι η \textbf{διαφορά} των $\mathrm{ΑΓ}$ και $\mathrm{ΑΒ}$: $\mathrm{ΒΓ}=\mathrm{ΑΓ}-\mathrm{ΑΒ}$.

\begin{center}
\begin{tikzpicture}
\draw[thick] (0,0)--(7,0);
\foreach \x/\l in {0/Α,2.5/Β,7/Γ}{\fill (\x,0) circle (2pt); \node[below] at (\x,-0.1) {$\mathrm{\l}$};}
\draw[gblue,thick] (0,0.4)--(2.5,0.4); \draw[gblue,thick,dashed] (2.5,0.4)--(7,0.4);
\node[gblue,above] at (1.25,0.45) {$\mathrm{ΑΒ}$}; \node[gblue,above] at (4.75,0.45) {$\mathrm{ΒΓ}$};
\end{tikzpicture}
\end{center}

\tmimasmall{Τεθλασμένη και περίμετρος}

Το μήκος μιας \textbf{τεθλασμένης γραμμής} είναι το άθροισμα των μηκών των ευθυγράμμων τμημάτων της. Το μήκος του ευθυγράμμου τμήματος $\mathrm{ΑΒ}$ είναι μικρότερο από το μήκος κάθε τεθλασμένης γραμμής με τα ίδια άκρα $\mathrm{Α}$, $\mathrm{Β}$ (εκτός αν η τεθλασμένη είναι το ίδιο το τμήμα). Δηλαδή το ευθύγραμμο τμήμα είναι ο συντομότερος δρόμος μεταξύ δύο σημείων.

\begin{center}
\begin{tikzpicture}
\draw[thick] (0,0)--(6,0); \draw[thick,gblue] (0,0)--(2,1.6)--(4,1)--(6,0);
\foreach \p/\l in {(0,0)/Α,(6,0)/Β,(2,1.6)/Γ,(4,1)/Δ}{\fill \p circle (2pt);}
\node[below left] at (0,0) {$\mathrm{Α}$}; \node[below right] at (6,0) {$\mathrm{Β}$}; \node[above] at (2,1.6) {$\Gamma$}; \node[above right] at (4,1) {$\Delta$};
\end{tikzpicture}
\end{center}

Το άθροισμα των πλευρών ενός ευθύγραμμου σχήματος (πολυγώνου) λέγεται \textbf{περίμετρος} του σχήματος.

\begin{prosoxi}
\begin{itemize}[leftmargin=1.2em,itemsep=0.3ex,topsep=2pt]
\item Η σχέση $\mathrm{ΑΓ}=\mathrm{ΑΒ}+\mathrm{ΒΓ}$ ισχύει μόνο όταν το $\mathrm{Β}$ βρίσκεται \emph{ανάμεσα} στα $\mathrm{Α}$ και $\mathrm{Γ}$ πάνω σε ευθεία. Αν δεν είναι συνευθειακά, το $\mathrm{ΑΓ}$ είναι μικρότερο από το $\mathrm{ΑΒ}+\mathrm{ΒΓ}$.
\item Η περίμετρος μετράει μόνο το μήκος του περιγράμματος, όχι την επιφάνεια.
\item Για την αφαίρεση προσέχω τη σειρά: $\mathrm{ΒΓ}=\mathrm{ΑΓ}-\mathrm{ΑΒ}$ και όχι $\mathrm{ΑΒ}-\mathrm{ΑΓ}$ (αρνητικό).
\item Για να συγκρίνω δύο διαδρομές πρέπει να έχουν την ίδια μονάδα μέτρησης.
\end{itemize}
\end{prosoxi}

\tmima{Λυμένα παραδείγματα}

\paradeigma{Τα δέντρα στη σειρά}
\tixero Τα δέντρα $\mathrm{Α},\,\mathrm{Β},\,\Gamma,\,\Delta$ είναι στη σειρά και $\mathrm{ΑΒ}=3$, $\mathrm{ΒΓ}=4$, $\Gamma\Delta=2{,}5$ μ.\par
\tithelo Τα $\mathrm{ΑΓ}$, $\mathrm{ΑΔ}$ και $\mathrm{ΒΔ}$.\par
\possindeo $\mathrm{ΑΓ}=\mathrm{ΑΒ}+\mathrm{ΒΓ}=3+4=7$. $\mathrm{ΑΔ}=\mathrm{ΑΓ}+\Gamma\Delta=7+2{,}5=9{,}5$. $\mathrm{ΒΔ}=\mathrm{ΑΔ}-\mathrm{ΑΒ}=9{,}5-3=6{,}5$.\par
\apantisi $7$, $9{,}5$ και $6{,}5$ μ.\par
\elegxos $\mathrm{ΒΔ}=\mathrm{ΒΓ}+\Gamma\Delta=4+2{,}5=6{,}5$.

\paradeigma{Με αφαίρεση}
\tixero Το $\mathrm{Β}$ είναι ανάμεσα στα $\mathrm{Α}$ και $\mathrm{Γ}$, $\mathrm{ΑΓ}=9$ και $\mathrm{ΑΒ}=4{,}5$.\par
\tithelo Το $\mathrm{ΒΓ}$.\par
\possindeo $\mathrm{ΒΓ}=\mathrm{ΑΓ}-\mathrm{ΑΒ}=9-4{,}5=4{,}5$.\par
\apantisi $4{,}5$.\par
\elegxos Άρα το $\mathrm{Β}$ είναι το μέσο του $\mathrm{ΑΓ}$.

\paradeigma{Ο κοντινός δρόμος}
\tixero Μονοπάτι $\mathrm{ΑΒΓ}$ με $\mathrm{ΑΒ}=6$ μ. και $\mathrm{ΒΓ}=8$ μ. (στρίβει στο $\mathrm{Β}$).\par
\tithelo Πόσο είναι το μήκος του και τι ξέρουμε για το $\mathrm{ΑΓ}$;\par
\possindeo Το μήκος της τεθλασμένης είναι $6+8=14$ μ. Το ευθύγραμμο τμήμα $\mathrm{ΑΓ}$ είναι συντομότερο.\par
\apantisi Τεθλασμένη $14$ μ. και $\mathrm{ΑΓ}<14$ μ.\par
\elegxos Αν το $\mathrm{Β}$ ήταν πάνω στο $\mathrm{ΑΓ}$, θα ήταν ίσα με $14$.

\paradeigma{Η περίμετρος}
\tixero Ένα ορθογώνιο παρτέρι με πλευρές $4{,}5$ μ. και $3{,}2$ μ.\par
\tithelo Το μήκος του φράχτη που το περιβάλλει.\par
\possindeo Η περίμετρος είναι το άθροισμα των πλευρών: $4{,}5+3{,}2+4{,}5+3{,}2=15{,}4$ μ.\par
\apantisi $15{,}4$ μ.\par
\elegxos $2\cdot(4{,}5+3{,}2)=2\cdot7{,}7=15{,}4$.

\paradeigma{Το μέσο και το άθροισμα}
\tixero Το $\mathrm{Γ}$ βρίσκεται πάνω στο $\mathrm{ΑΒ}$ με $\mathrm{ΑΓ}=3$ και $\mathrm{ΓΒ}=7$ και το $\mathrm{Μ}$ είναι το μέσο του $\mathrm{ΑΒ}$.\par
\tithelo Το $\mathrm{ΓΜ}$.\par
\possindeo $\mathrm{ΑΒ}=3+7=10$, άρα $\mathrm{ΑΜ}=5$. Το $\mathrm{Μ}$ είναι πέρα από το $\mathrm{Γ}$ και $\mathrm{ΓΜ}=\mathrm{ΑΜ}-\mathrm{ΑΓ}=5-3=2$.\par
\apantisi $\mathrm{ΓΜ}=2$.\par
\elegxos $\mathrm{ΓΜ}+\mathrm{ΜΒ}=2+5=7=\mathrm{ΓΒ}$.

\tmima{Ασκήσεις}

\zoni{Το κατάλαβα}{άθροισμα, διαφορά, περίμετρος}

\ex Τα $\mathrm{Α},\,\mathrm{Β},\,\Gamma,\,\Delta$ είναι στη σειρά σε μια ευθεία με $\mathrm{ΑΒ}=3$, $\mathrm{ΒΓ}=4$, $\Gamma\Delta=2{,}5$. Βρες τα $\mathrm{ΑΓ}$, $\mathrm{ΑΔ}$, $\mathrm{ΒΔ}$.

\ex Το $\mathrm{Β}$ είναι ανάμεσα στα $\mathrm{Α}$ και $\mathrm{Γ}$. Αν $\mathrm{ΑΓ}=9$ και $\mathrm{ΑΒ}=4{,}5$, βρες το $\mathrm{ΒΓ}$.

\ex Βρες το μήκος της τεθλασμένης $\mathrm{ΑΒΓΔ}$ με $\mathrm{ΑΒ}=5$, $\mathrm{ΒΓ}=3{,}5$, $\Gamma\Delta=6$.

\ex Βρες την περίμετρο: α) ορθογωνίου $4{,}5$ και $3{,}2$ \quad β) τριγώνου με πλευρές $5$, $6$, $7{,}5$ \quad γ) τετραγώνου πλευράς $2{,}8$ \quad δ) κανονικού εξαγώνου πλευράς $4$.

\ex Στην ευθεία $\mathrm{Α},\,\mathrm{Β},\,\Gamma$ με $\mathrm{ΑΒ}=2\mathrm{ΒΓ}$ και $\mathrm{ΑΓ}=12$ βρες τα $\mathrm{ΑΒ}$ και $\mathrm{ΒΓ}$.

\ex Σωστό ή λάθος; α) Αν το $\mathrm{Β}$ είναι ανάμεσα στα $\mathrm{Α}$ και $\Gamma$ τότε $\mathrm{ΑΓ}=\mathrm{ΑΒ}+\mathrm{ΒΓ}$. \quad β) Η τεθλασμένη μεταξύ δύο σημείων είναι πάντα μεγαλύτερη από το τμήμα που τα ενώνει (ή ίση). \quad γ) Η περίμετρος ενός τετραγώνου πλευράς $\alpha$ είναι $\alpha^2$. \quad δ) $\mathrm{ΒΓ}=\mathrm{ΑΒ}-\mathrm{ΑΓ}$ όταν το $\mathrm{Β}$ είναι ανάμεσα στα $\mathrm{Α}$, $\mathrm{Γ}$.

\zoni{Το χρησιμοποιώ}{φράχτες, δρόμοι, κήποι}

\ex \textbf{Ο φράχτης.} Ένα ορθογώνιο παρτέρι $4{,}5$ μ. επί $3{,}2$ μ. θα περιφραχθεί με δίχτυ που κοστίζει $3$ € το μέτρο. Πόσο θα κοστίσει; Μην ξεχάσεις ότι ένα μέτρο το χρειάζεται για την είσοδο που μένει ανοιχτή.

\ex \textbf{Ο κοντινός δρόμος.} Το μονοπάτι $\mathrm{ΑΒΓ}$ έχει $\mathrm{ΑΒ}=6$ μ., $\mathrm{ΒΓ}=8$ μ. Ο Νίκος κόβει δρόμο πηγαίνοντας ευθεία από το $\mathrm{Α}$ στο $\mathrm{Γ}$. Έκανε περισσότερο ή λιγότερο δρόμο από $14$ μ.; Γιατί;

\ex \textbf{Η απόσταση.} Τρία σπίτια $\mathrm{Α}$, $\mathrm{Β}$, $\Gamma$ βρίσκονται στον ίδιο δρόμο, με το $\mathrm{Β}$ στη μέση. Το $\mathrm{Α}$ απέχει από το $\mathrm{Β}$ $350$ μ. και το $\mathrm{Β}$ από το $\Gamma$ $480$ μ. Πόσο απέχουν τα $\mathrm{Α}$ και $\Gamma$;

\ex \textbf{Η τεθλασμένη διαδρομή.} Η Πολυξένη πηγαίνει από το σχολείο στο σπίτι με στάσεις: $250$ μ., $180$ μ. και $320$ μ. Αν πήγαινε ευθεία θα έκανε $600$ μ. Πόσο περισσότερο περπάτησε;

\ex \textbf{Ο φράχτης του κήπου.} Ο κήπος έχει σχήμα πενταγώνου με πλευρές $12$, $9$, $14$, $10$ και $8$ μ. Πόσος φράχτης χρειάζεται;

\zoni{Το σκέφτομαι}{ο ευθύς είναι ο συντομότερος}

\ex Εξήγησε γιατί το ευθύγραμμο τμήμα είναι ο συντομότερος δρόμος ανάμεσα σε δύο σημεία (σκέψου ένα τεντωμένο νήμα).

\ex Αν $\mathrm{ΑΒ}=5$ και $\mathrm{ΒΓ}=7$ και τα $\mathrm{Α},\,\mathrm{Β},\,\Gamma$ δεν είναι συνευθειακά, ανάμεσα σε ποιες τιμές βρίσκεται το $\mathrm{ΑΓ}$;

\ex Ένα σχοινί $\mathrm{ΑΒ}$ μήκους $10$ μ. διπλώνεται ώστε το $\mathrm{Α}$ να πέσει πάνω στο $\mathrm{Β}$. Πού βρίσκεται η πτυχή;

\ex Δείξε με παράδειγμα ότι η περίμετρος δύο σχημάτων μπορεί να είναι ίση χωρίς τα σχήματα να είναι ίσα.

\zoni{Για το σπίτι}{μέτρησε και πρόσθεσε}

\ex Μέτρησε τις πλευρές του δωματίου σου και υπολόγισε την περίμετρο της κάτοψής του.

\ex Φτιάξε μια τεθλασμένη με τρία τμήματα από σπάγκο και σύγκρινε το μήκος της με την απόσταση των άκρων.
